\documentclass[11pt,a4paper]{article}
\usepackage[margin=2.4cm]{geometry}
\usepackage{amsmath,amssymb,amsthm}
\usepackage[dvipsnames]{xcolor}
\usepackage{listings}
\usepackage[colorlinks=true,linkcolor=blue!50!black,urlcolor=blue!50!black]{hyperref}

\newcommand{\labeledmatrixTwo}[2]{
    \begin{array}{c@{}c@{}c}
        & \textcolor{gray}{\begin{array}{cc}
            0 & 1
        \end{array}} \\[5pt]
        \textcolor{gray}{\begin{array}{c}
            0 \\ 1
        \end{array}} &
        \left[ \begin{array}{cc}
            #1 \\ #2 \\
        \end{array} \right]
    \end{array}
}

% Define a macro for the labeled 4x4 matrix with gray labels, with compact content columns
\newcommand{\labeledmatrixFour}[4]{
    \begin{array}{c@{}c}
        & \textcolor{gray}{\begin{array}{cccc}
            00 & 01 & 10 & 11
        \end{array}} \\[5pt]
        \textcolor{gray}{\begin{array}{c}
            00 \\ 01 \\ 10 \\ 11
        \end{array}} &
        \left[ \begin{array}{@{\hspace{12pt}}c@{\hspace{12pt}}c@{\hspace{12pt}}c@{\hspace{12pt}}c}
            #1 \\ #2 \\ #3 \\ #4 \\
        \end{array} \right]
    \end{array}
}

% Define a macro for the labeled 8x8 matrix with gray labels, with compact content columns
\newcommand{\labeledmatrixEight}[8]{
    \begin{array}{c@{}c}
        & \textcolor{gray}{\begin{array}{cccccccc}
            000 & 001 & 010 & 011 & 100 & 101 & 110 & 111
        \end{array}} \\[5pt]
        \textcolor{gray}{\begin{array}{c}
            000 \\ 001 \\ 010 \\ 011 \\ 100 \\ 101 \\ 110 \\ 111
        \end{array}} &
        \left[ \begin{array}{@{\hspace{15pt}}c@{\hspace{15pt}}c@{\hspace{15pt}}c@{\hspace{15pt}}c@{\hspace{15pt}}c@{\hspace{15pt}}c@{\hspace{15pt}}c@{\hspace{15pt}}c}
            #1 \\ #2 \\ #3 \\ #4 \\ #5 \\ #6 \\ #7 \\ #8 \\
        \end{array} \right]
    \end{array}
}

\newcommand{\ket}[1]{\ensuremath{\left| #1 \right\rangle}}
\newcommand{\bra}[1]{\ensuremath{\left\langle #1 \right|}}
\newcommand{\sqrtf}[2]{\sqrt{\frac{#1}{#2}}}
\newcommand{\norm}[1]{\left\| #1 \right\|}
\newcommand{\braket}[2]{\left\langle #1 \middle| #2 \right\rangle}
\newcommand{\ketbra}[2]{\left| #1 \middle\rangle \middle\langle #2 \right|}
\newcommand{\innerw}[3]{\left\langle #1 \middle| #2 \middle| #3 \right\rangle}
\newcommand{\inner}[2]{\left\langle #1 \middle| #2 \right\rangle}
\newcommand{\setof}[2]{\left\{ #1 \middle| #2 \right\}}
\newcommand{\bigO}[1]{\mathcal{O}\left( #1 \right)}
\newcommand{\Pauli}{\mathcal{P}}
\newcommand{\optype}[1]{\mathsf{#1}}
\newcommand{\mattype}[1]{\mathbf{#1}}
\newcommand{\vecttype}[1]{\mathbf{#1}}
\newcommand{\XX}{\mattype{X}}
\newcommand{\YY}{\mattype{Y}}
\newcommand{\ZZ}{\mattype{Z}}  
\newcommand{\II}{\mattype{I}}
\newcommand{\UU}{\mattype{U}}
\newcommand{\RR}{\mattype{R}}
\newcommand{\nn}{\vecttype{n}}
\newcommand{\xx}{\vecttype{x}}
\newcommand{\zz}{\vecttype{z}}
\newcommand{\aas}{\vecttype{a}}
\newcommand{\sss}{\vecttype{s}}

\newcommand{\R}{\mathbb{R}}
\newcommand{\C}{\mathbb{C}}
\newcommand{\Z}{\mathbb{Z}}
\newcommand{\Stabgroup}[1]{\mathcal{S}\left( #1 \right)}
\newcommand{\Stab}[1]{\Stabgroup{#1}}
\newcommand{\Nstab}[1]{\mathcal{N}\left( #1 \right)}
\newcommand{\mat}[1]{\begin{bmatrix}#1\end{bmatrix}}
\newcommand{\vect}[1]{\left[\begin{array}{r} #1 \end{array}\right]}
\newcommand{\boldon}[2]{\alt<#1>{\bm{#2}}{#2}}


\newcommand{\measure}[2]{\textsf{measure}_{#1, #2}\,}
\newcommand{\hilbert}{\mathfrak{h}}
\newcommand{\Hilbert}{\mathcal{H}}


% Colors
\definecolor{codebg}{rgb}{0.98,0.98,0.98}
\definecolor{codekw}{rgb}{0.10,0.10,0.60}
\definecolor{codecmt}{rgb}{0.00,0.45,0.00}
\definecolor{codestr}{rgb}{0.55,0.00,0.00}

\lstset{
  language=Haskell,
  backgroundcolor=\color{codebg},
  basicstyle=\ttfamily\small,
  keywordstyle=\color{codekw}\bfseries,
  commentstyle=\color{codecmt}\itshape,
  stringstyle=\color{codestr},
  showstringspaces=false,
  numbers=left,
  numberstyle=\tiny\color{gray},
  stepnumber=1,
  numbersep=8pt,
  frame=single,
  framerule=0.2pt,
  tabsize=2,
  breaklines=true,
  breakatwhitespace=true,
  columns=fullflexible,
  keepspaces=true,
}

% optional: nicer Unicode/math-like symbols in code
\lstset{
  literate=
    {→}{{$\to$}}1
    {←}{{$\leftarrow$}}1
    {⇒}{{$\Rightarrow$}}1
    {∀}{{$\forall$}}1
    {∃}{{$\exists$}}1
}
\lstset{literate=
    {→}{{$\to$}}1 {←}{{$\leftarrow$}}1 {⇒}{{$\Rightarrow$}}1
    {∀}{{$\forall$}}1 {∃}{{$\exists$}}1
    {⊗}{{$\otimes$}}1 {∘}{{$\circ$}}1 {⊕}{{$\oplus$}}1
    {ψ}{{$\psi$}}1 {φ}{{$\varphi$}}1 {θ}{{$\theta$}}1 {≈}{{$\approx$}}1 {π}{{$\pi$}}1
    {⟨}{{$\langle$}}1 {⟩}{{$\rangle$}}1 {·}{{$\cdot$}}1
    {²}{{$^2$}}1 {±}{{$\pm$}}1
}

\theoremstyle{definition}
\newtheorem{exercise}{Exercise}[section]
% Complete, executable definitions use qppcode.  The checker extracts these
% environments from this file; ordinary lstlisting blocks are transcripts,
% declarations quoted from QPP, or deliberately incomplete exercises.
\lstnewenvironment{qppcode}{\lstset{language=Haskell}}{}

\title{Introduction to Quantum Computing\\
       {\large with the Quantum Programming Playground}\\[1ex]
       {\normalsize Lecture 0 notes --- ATPL 2025}}
\author{James Avery (\texttt{avery@di.ku.dk})}
\date{November 2025}

\begin{document}
\maketitle

These notes introduce the mathematical model of quantum computation and its
implementation in QPP, the \emph{Quantum Programming Playground}.  They are
written for two audiences.  Computer-science students already know Haskell
but may be meeting quantum computing for the first time; quantum-information
students know the mathematics but may be new to functional programming.
The first week's work connects the two subjects through QPP.

We begin with three principles:
\[
  \text{states are unit vectors},\qquad
  \text{gates are unitary maps},\qquad
  \text{measurement follows the Born rule}.
\]
Each idea appears as mathematics and as a QPP value.  In the exercises, derive
a prediction, implement it, and explain the result.  Haskell programmers
should slow down at the derivations; quantum-information students should pay
particular attention to types and recursion.

\paragraph{Starting QPP.}
From the QPP repository, run \lstinline|cabal repl lib:qpp| and import:
\begin{lstlisting}
import QPP
import QPP.Semantics.Matrix
import QPP.PrettyPrint
import Data.Complex
\end{lstlisting}
Put definitions in a file such as \texttt{Playground.hs}; load it with
\lstinline|:add Playground.hs| and reload after edits with \lstinline|:r|.
Predict the result of each example before evaluating it.

\section{Enough Haskell to begin}
\label{sec:haskell}

This section is a quick orientation.
If the language is familiar, concentrate on the QPP types and notation.  If
it is new, enter the examples in GHCi and keep your definitions in
\texttt{Playground.hs}.  Type errors often reveal a mismatch between the
program and the mathematics.

\subsection{Expressions, types, and functions}

GHCi evaluates expressions.  The command \lstinline|:type| (abbreviated
\lstinline|:t|) reports an expression's type without evaluating it:
\begin{lstlisting}
ghci> 2 ^ 10
1024
ghci> :t sqrt
sqrt :: Floating a => a -> a
ghci> :t [1,2,3]
[1,2,3] :: Num a => [a]
\end{lstlisting}
The constraint \lstinline|Floating a =>| permits any type \lstinline|a| that
supports floating-point operations, not only \lstinline|Double|.
A signature \lstinline|f :: A -> B| says that \lstinline|f| takes an
\lstinline|A| and returns a \lstinline|B|.  Function application needs no
parentheses: \lstinline|sqrt 2| means $\sqrt2$.  Application binds more
tightly than infix operators, so \lstinline|magnitude a ^ 2| means
\lstinline|(magnitude a) ^ 2|.

Definitions are equations, and names denote values rather than mutable
variables.  This becomes important when randomness is made an explicit input
in Section~\ref{sec:measurement}.

\subsection{Complex numbers and lists}

Quantum amplitudes are complex numbers.  QPP uses
\lstinline|ComplexT| as an abbreviation for \lstinline|Complex Double|.
The constructor \lstinline|a :+ b| denotes $a+bi$:
\begin{lstlisting}
ghci> magnitude (0.6 :+ 0.8)
1.0
ghci> conjugate (2 :+ 3)
2.0 :+ (-3.0)
\end{lstlisting}
A bare numeric literal can be used as a complex number when the surrounding
type requires it; for example, \lstinline|[3,4] :: [ComplexT]| denotes two
real amplitudes.

Lists occur throughout QPP.  A list comprehension transforms or enumerates
their elements:
\begin{lstlisting}
ghci> [x * x | x <- [1,2,3,4]]
[1,4,9,16]
ghci> [(a,b) | a <- [0,1], b <- [0,1]]
[(0,0),(0,1),(1,0),(1,1)]
\end{lstlisting}
The second form will later enumerate the basis states of two qubits.
Ordinary library functions such as \lstinline|map|, \lstinline|sum|, and
\lstinline|length| are equally useful.

As a first example, the Euclidean norm of a list of amplitudes is:
\begin{qppcode}
norm2 :: [ComplexT] -> Double
norm2 as = sqrt (sum [magnitude a ^ 2 | a <- as])

ghci> norm2 [3,4]
5.0
\end{qppcode}

\begin{exercise}[Types as specifications]
Predict the types of \lstinline|magnitude|, \lstinline|sum|,
\lstinline|map magnitude|, and \lstinline|[0.6 :+ 0.8, 1]|; then ask GHCi.
Some inferred types will be more general than expected.  Read the type-class
constraints as requirements on the operations a type must support.

Then implement:
\begin{lstlisting}
average :: [Double] -> Double
timesI  :: ComplexT -> ComplexT
\end{lstlisting}
Use \lstinline|fromIntegral| to convert the result of \lstinline|length|.
Check that \lstinline|average [1,2,3]| is $2$ and that applying
\lstinline|timesI| twice multiplies a number by $-1$.
\end{exercise}

\section{States, amplitudes, and probabilities}
\label{sec:states}

\subsection{A qubit is a vector}

A classical bit has one of two values.  A qubit has a state in the
two-dimensional complex vector space $\C^2$.  We name its standard basis
$\ket0$ and $\ket1$, so every one-qubit state has the form
\[
  \ket\psi = a\ket0+b\ket1,
  \qquad a,b\in\C.
\]
The coefficients $a$ and $b$ are \emph{amplitudes}.  We represent a pure state
by a normalised vector:
\[
  \norm{\psi}^2=|a|^2+|b|^2=1.
\]
For example,
\[
  \ket\psi=\sqrt{\tfrac13}\ket0+\sqrt{\tfrac23}\ket1
\]
is a valid state.  In QPP:
\begin{qppcode}
psi :: StateT
psi = sqrt (1/3) .* ket [0] .+ sqrt (2/3) .* ket [1]

ghci> printS psi
0.5773502691896257*(ket [0]) + 0.816496580927726*(ket [1])
ghci> norm psi
1.0
\end{qppcode}
Here \lstinline|ket [0]| constructs $\ket0$, \lstinline|.*| is scalar
multiplication, and \lstinline|.+| is vector addition.  QPP uses
floating-point complex numbers, so small rounding errors are normal.
\lstinline|StateT| is the matrix backend's state representation; use these
operations rather than relying on its internals.

\subsection{The Born rule}

The amplitudes are not directly observable.  Measuring the state above in
the computational basis returns a classical bit.  The Born rule assigns the
probabilities
\[
  \Pr(0)=|a|^2,\qquad \Pr(1)=|b|^2.
\]
Thus measuring $\ket\psi$ gives $0$ with probability $1/3$ and $1$ with
probability $2/3$.

More generally, the inner product is
\[
  \inner{\varphi}{\psi}=\sum_k \overline{\varphi_k}\psi_k,
\]
and the probability of outcome $\varphi$ in an orthonormal measurement basis,
when the system is in the normalised state $\ket\psi$, is
\[
  \Pr(\varphi\mid\psi)=|\inner{\varphi}{\psi}|^2.
\]
This gives a compact QPP definition:
\begin{qppcode}
prob :: StateT -> StateT -> Double
prob φ ψ = magnitude (inner φ ψ) ^ 2

ghci> [prob (ket [k]) psi | k <- [0,1]]
[0.3333333333333333,0.6666666666666666]
\end{qppcode}
The function \lstinline|prob| calculates a probability; it does not perform a
measurement or change the state.  Section~\ref{sec:measurement} introduces
\lstinline|Measure|, which samples an outcome and updates the state.
The conjugation in the first argument matters.  For example,
$\inner{i\psi}{\psi}=-i\inner{\psi}{\psi}$.

The type of \lstinline|inner| also illustrates a central Haskell convention:
\begin{lstlisting}
inner :: StateT -> StateT -> ComplexT
\end{lstlisting}
Arrows associate to the right, so this means
\lstinline|StateT -> (StateT -> ComplexT)|.  Supplying only the first
argument produces a new function:
\begin{qppcode}
bra :: StateT -> (StateT -> ComplexT)
bra φ = inner φ

ghci> bra (ket [1]) psi
0.816496580927726 :+ 0.0
\end{qppcode}
This is \emph{partial application}.  Dirac notation makes the same split:
$\bra\varphi$ maps $\ket\psi$ to $\inner{\varphi}{\psi}$.  The notation and
the function type have the same shape.

\begin{exercise}[Amplitudes and probabilities]
Define
\begin{lstlisting}
born :: [ComplexT] -> [Double]
\end{lstlisting}
so that \lstinline|born as| returns $|a_k|^2$ for every amplitude in
\lstinline|as|.  Then define
\begin{lstlisting}
normalizeAmps :: [ComplexT] -> [ComplexT]
\end{lstlisting}
which normalises a nonzero amplitude vector.  Check the result on
\lstinline|[3,4]|.  Check also that the entries returned by \lstinline|born|
sum to approximately $1$ after normalisation.  What should the normaliser do
on the zero vector?
\end{exercise}

\subsection{A state depends on the chosen basis}

The computational basis is convenient, but not privileged.  Define
\[
  \ket+=\frac{\ket0+\ket1}{\sqrt2},\qquad
  \ket-=\frac{\ket0-\ket1}{\sqrt2}.
\]
These vectors are orthonormal and therefore form another basis of $\C^2$.
In QPP they are obtained by applying $H$ to the computational basis:
\begin{qppcode}
plus  = apply (evalOp H) (ket [0])
minus = apply (evalOp H) (ket [1])
\end{qppcode}
They have identical computational-basis probabilities, but their relative
sign can affect a later computation.  To measure in the $\pm$ basis, apply
$H$ and then measure in the computational basis.  Indeed,
$H\ket+=\ket0$ and $H\ket-=\ket1$.

Multiplying an entire state by a phase $e^{i\delta}$ changes no measurement
probability, because
\[
  |\inner{\varphi}{e^{i\delta}\psi}|^2
  =|e^{i\delta}|^2|\inner{\varphi}{\psi}|^2
  =|\inner{\varphi}{\psi}|^2.
\]
Such a \emph{global phase} is unobservable.  A phase between two amplitudes
is a \emph{relative phase}; it can be observed through interference.

\begin{exercise}[Relative phase becomes an outcome]
Construct
\begin{qppcode}
wPlus  = normalize (3 .* ket [0] .+ 4 .* ket [1])
wMinus = normalize (3 .* ket [0] .- 4 .* ket [1])
\end{qppcode}
First derive and then compute their computational-basis probabilities.  Next
apply $H$ to both states and repeat.  Why can the second experiment distinguish
the states when the first cannot?  Express the calculation with one Haskell
function that accepts a state and returns its two probabilities.
\end{exercise}

\section{Quantum gates and QPP}
\label{sec:gates}

\subsection{Unitary evolution}

Between measurements an isolated quantum system evolves linearly without
changing norms.  In finite dimensions this means that evolution is given by
a unitary matrix $U$, satisfying
\[
  U^\dagger U=I.
\]
Unitarity has two immediate consequences: the output is still a valid state,
and the evolution is reversible, with inverse $U^\dagger$.

Four basic one-qubit gates are
\[
  X=\begin{pmatrix}0&1\\1&0\end{pmatrix},\quad
  Y=\begin{pmatrix}0&-i\\i&0\end{pmatrix},\quad
  Z=\begin{pmatrix}1&0\\0&-1\end{pmatrix},\quad
  H=\frac1{\sqrt2}\begin{pmatrix}1&1\\1&-1\end{pmatrix}.
\]
$X$ exchanges $\ket0$ and $\ket1$,
$Y\ket0=i\ket1$ and $Y\ket1=-i\ket0$, while $Z$ changes the sign of the
$\ket1$ amplitude.  Finally, $H$ changes between the computational and
$\pm$ bases:
\[
  H\ket0=\ket+,\quad H\ket1=\ket-,\qquad
  H\ket+=\ket0,\quad H\ket-=\ket1.
\]

\subsection{Operators as syntax trees}

QPP represents an operator by a value of the algebraic data type \lstinline|QOp|.
The relevant constructors are:
\begin{lstlisting}
data QOp = Id Nat | X | Y | Z | H | SX
         | R QOp Rational
         | C QOp
         | Tensor QOp QOp
         | Compose QOp QOp
         | Adjoint QOp
         -- further constructors omitted
\end{lstlisting}
Here \lstinline|I| abbreviates \lstinline|Id 1|, the one-qubit identity, and
\lstinline|SX| is a square root of $X$: \lstinline|SX <> SX| has the same
matrix as \lstinline|X|.

Constructing a \lstinline|QOp| builds a syntax tree; it does not yet multiply
matrices.  This separation lets us inspect and transform operator syntax before
choosing how to execute it.
\begin{lstlisting}
ghci> H <> H
Compose H H
ghci> C X
C X
\end{lstlisting}
For readers new to Haskell, each line of the declaration introduces a
\emph{constructor}.  A constructor is a function for building values:
\lstinline|R| stores an axis and an angle, \lstinline|C| stores a
sub-operator, and \lstinline|Tensor| stores two sub-operators.  Consequently,
a value such as \lstinline|C (H ⊗ X)|
is a tree whose root is \lstinline|C| and whose child is a
\lstinline|Tensor| node.  This is why QPP can transform an operator expression
without evaluating it.

Pattern matching takes such values apart.  For example, a small analysis can
recognise the outermost constructor:
\begin{qppcode}
hasOuterControl :: QOp -> Bool
hasOuterControl (C _) = True
hasOuterControl _     = False
\end{qppcode}
The first pattern matches a \lstinline|C| node and ignores its child; the
second matches everything else.  Recursive analyses also inspect children.

The matrix backend gives that syntax a meaning:
\begin{center}
\begin{tabular}{ll}
\lstinline|evalOp u| & evaluate an operator as a matrix \\
\lstinline|apply m s| & apply a matrix to a state \\
\lstinline|printM m| & display a matrix \\
\lstinline|printS s| & display a state
\end{tabular}
\end{center}
For example:
\begin{lstlisting}
ghci> printM (evalOp X)
 [[0,1],
  [1,0]]
ghci> printS (apply (evalOp H) (ket [0]))
0.7071067811865475*(ket [0]) + 0.7071067811865475*(ket [1])
\end{lstlisting}

Composition is written \lstinline|u <> v| or \lstinline|u ∘ v| and follows
the mathematical convention: apply the right-hand operator first.  Tensor
product, written \lstinline|u ⊗ v| (ASCII: \lstinline|u <.> v|), places
operators on separate registers.  Finally, \lstinline|adj u| constructs the
adjoint:
\begin{lstlisting}
ghci> printS (apply (evalOp (adj SX <> SX)) (ket [0]))
1*(ket [0])
\end{lstlisting}

\begin{exercise}[Reading operator syntax]
Predict the result of each expression before running it:
\begin{lstlisting}
apply (evalOp (X ⊗ I)) (ket [0,1])
apply (evalOp (H <> H)) (ket [0])
apply (evalOp (H <> Z <> H)) (ket [0])
\end{lstlisting}
Explain the third result by following the state through the three gates.
\end{exercise}

\begin{exercise}[Structural recursion]
Write \lstinline|arity :: QOp -> Int| by pattern matching on the constructors
above.  A control adds one qubit; a tensor product adds arities; an adjoint
and a rotation preserve arity.  Assume compositions are well formed, so their
operands have equal arity.
Start with:
\begin{lstlisting}
arity :: QOp -> Int
arity (Id n)       = n
arity X            = 1
arity (C u)        = 1 + arity u
arity (Tensor u v) = arity u + arity v
-- complete the remaining cases
\end{lstlisting}
The patterns identify the outer constructor and name its fields; calls on
\lstinline|u| and \lstinline|v| recurse on smaller trees.  Test against the
library on every constructor, including a rotation about a multi-qubit
operator.

($\star$) Use \lstinline|Maybe| to make composition failure explicit.  Assuming
the children are otherwise valid, its composition case is:
\begin{lstlisting}
checkedArity (Compose u v) =
  case (checkedArity u, checkedArity v) of
    (Just m, Just n) | m == n -> Just m
    _                         -> Nothing
\end{lstlisting}
This separates validation from the structural calculation; the exercise sheet
extends the idea to a complete validator.
\end{exercise}

\section{Interference}
\label{sec:interference}

Apply $H$ twice to $\ket0$:
\[
  \ket0 \xrightarrow{H} \frac{\ket0+\ket1}{\sqrt2}
  \xrightarrow{H}
  \frac{(\ket0+\ket1)+(\ket0-\ket1)}2=\ket0.
\]
The two contributions to the $\ket1$ amplitude cancel.  This is
\emph{interference}.  It is misleading to imagine that the first $H$ tosses
a coin: no random event occurs until measurement.  Amplitudes evolve and may
cancel; classical probabilities cannot.

A $Z$-rotation supplies a controllable relative phase:
\[
  R_Z(\theta)=
  \begin{pmatrix}
    e^{-i\theta\pi/2}&0\\0&e^{i\theta\pi/2}
  \end{pmatrix}.
\]
QPP measures the parameter in units of $\pi$, so this is \lstinline|R Z θ|.
Placing the rotation between two Hadamard gates gives
\[
  \ket0 \xrightarrow{H} \ket+
  \xrightarrow{R_Z(\theta)}
  \frac{e^{-i\theta\pi/2}\ket0+e^{i\theta\pi/2}\ket1}{\sqrt2}
  \xrightarrow{H}
  \cos\!\left(\frac{\theta\pi}{2}\right)\ket0
  -i\sin\!\left(\frac{\theta\pi}{2}\right)\ket1.
\]
Hence
\[
  \Pr(0)=\cos^2\!\left(\frac{\theta\pi}{2}\right),\qquad
  \Pr(1)=\sin^2\!\left(\frac{\theta\pi}{2}\right).
\]
This circuit is the qubit model of a Mach--Zehnder interferometer: the two
Hadamards act as beam splitters and the $Z$-rotation changes one path's
phase.
\begin{qppcode}
mziOp :: Rational -> QOp
mziOp θ = H ∘ R Z θ ∘ H

ghci> prob (ket [0]) (apply (evalOp (mziOp (1/4))) (ket [0]))
0.8535533905932733
\end{qppcode}

\begin{exercise}[The interference fringe]
Compute $\Pr(0)$ at $\theta=0,\frac14,\frac12,\frac34,1$ on paper and
define
\begin{lstlisting}
p0 :: Rational -> Double
\end{lstlisting}
to compute the same values from \lstinline|mziOp| using \lstinline|prob|.
Generate the five results with \lstinline|map| or a comprehension and compare
them with the formula.  Explain the endpoints using constructive and
destructive interference.
\end{exercise}

\section{Many qubits and entanglement}
\label{sec:many}

The state space of two qubits is the tensor product
\[
  \C^2\otimes\C^2\cong\C^4,
\]
with basis $\ket{00},\ket{01},\ket{10},\ket{11}$.  In general, an
$n$-qubit state is
\[
  \ket\psi=\sum_{x\in\{0,1\}^n}a_x\ket x,
  \qquad \sum_x|a_x|^2=1.
\]
It therefore has $2^n$ amplitudes.  Quantum algorithms manipulate these
amplitudes through structured operations and interference; they cannot simply
read them out, since one measurement returns only classical outcomes.

Some joint states factor into states of their parts.  For example,
\[
  \ket+\otimes\ket0=\frac{\ket{00}+\ket{10}}{\sqrt2}
\]
is a \emph{product state}.  A state that does not factor is
\emph{entangled}.  The standard example is the Bell state
\[
  \ket{\Phi^+}=\frac{\ket{00}+\ket{11}}{\sqrt2}.
\]

The controlled-$X$ gate, \lstinline|C X|, flips its target precisely when
its control is $1$.  It produces the Bell state from $\ket{00}$:
\begin{qppcode}
bellOp :: QOp
bellOp = C X ∘ (H ⊗ I)

bell :: StateT
bell = apply (evalOp bellOp) (ket [0,0])

ghci> printS bell
0.7071067811865475*(ket [0,0]) + 0.7071067811865475*(ket [1,1])
ghci> [prob (ket [a,b]) bell | a <- [0,1], b <- [0,1]]
[0.4999999999999999,0.0,0.0,0.4999999999999999]
\end{qppcode}
Measuring either qubit in the computational basis gives an unbiased bit, and
the outcomes are perfectly correlated.  A shared classical random bit could
produce the same distribution.  What makes the Bell state entangled is that
its amplitude vector cannot be written as a tensor product of one-qubit state
vectors.

\begin{exercise}[Why the Bell state is entangled]
Expand
\[
  (a_0\ket0+a_1\ket1)\otimes(b_0\ket0+b_1\ket1).
\]
Assume it equals $\ket{\Phi^+}$ and use the four coefficient equations to
derive a contradiction.  A nonzero two-qubit amplitude vector is a product
exactly when
\[
  a_{00}a_{11}-a_{01}a_{10}=0.
\]
Derive this condition, then implement it for a four-element amplitude list:
\begin{lstlisting}
productDet :: [ComplexT] -> ComplexT
productDet [a00,a01,a10,a11] = undefined -- complete this expression
productDet _                  = error "expected four amplitudes"
\end{lstlisting}
Evaluate it on the Bell amplitudes and on
$\frac12\bigl((\ket0+\ket1)\otimes(\ket0-\ket1)\bigr)$.
\end{exercise}

\begin{exercise}[Four Bell states]
Starting from \lstinline|bellOp|, produce the other three Bell states
\[
 \frac{\ket{00}-\ket{11}}{\sqrt2},\quad
 \frac{\ket{01}+\ket{10}}{\sqrt2},\quad
 \frac{\ket{01}-\ket{10}}{\sqrt2}
\]
by applying $X$ and $Z$ to one qubit.  A result that differs only by global
phase represents the same physical state.  Write a function that returns all
four computational-basis probabilities of a two-qubit state, and use it to
compare the four Bell states.  Do those probabilities distinguish them?
\end{exercise}

\begin{exercise}[Bell-basis measurement]
Apply \lstinline|adj bellOp| to each Bell state.  Derive which computational
basis ket each state becomes, then check all four mappings in QPP.  After
Section~\ref{sec:measurement}, append a measurement to obtain a Bell-basis
measurement program.
\end{exercise}

\section{Programs with measurement}
\label{sec:measurement}

Measurement is probabilistic and changes the state.  When a computational-
basis measurement returns $x$, the state collapses to the component
consistent with $x$, renormalised.  QPP distinguishes unitary operators from
programs that may measure:
\begin{lstlisting}
data Step = Unitary QOp
          | Initialize [Nat] [Bool] -- reset selected qubits
          | Measure [Nat]

type Program = [Step]
\end{lstlisting}
Unlike the \lstinline|data| declaration, the \lstinline|type| declaration
does not introduce a new representation: it gives \lstinline|[Step]| a useful
name.  Sequential program structure is therefore represented by an ordinary
Haskell list, while each element records what kind of step to perform.
For example, the following program prepares $\ket+$ and measures it:
\begin{qppcode}
coin :: Program
coin = [Unitary H, Measure [0]]
\end{qppcode}

The interpreter receives its random choices as a list of numbers in $[0,1)$.
Its interface and result types are:
\begin{lstlisting}
type RNG      = [Double]
type Outcomes = [Bool] -- False = 0, True = 1

evalProg :: Program -> StateT -> RNG
         -> (StateT, Outcomes, RNG)
\end{lstlisting}
The expression \lstinline|repeat x| is the infinite list
\lstinline|[x,x,...]|.  Such constant streams let us force either side of a
fair measurement:
\begin{qppcode}
choose0, choose1 :: RNG
choose0 = repeat 0.0
choose1 = repeat 0.999

ghci> take 3 choose0
[0.0,0.0,0.0]
ghci> let (_, o0, _) = evalProg coin (ket [0]) choose0
ghci> let (_, o1, _) = evalProg coin (ket [0]) choose1
ghci> (o0, o1)
([False],[True])
\end{qppcode}
Haskell evaluates only the prefix demanded by \lstinline|take| or
\lstinline|evalProg|, so a
pure function can receive an unbounded supply of choices and return the unused
tail.  The patterns above deconstruct result triples; \lstinline|_| ignores
components we do not need.

For sampling, replace the test streams by a seeded pseudorandom stream:
\begin{lstlisting}
import System.Random (mkStdGen, randoms)
\end{lstlisting}
\begin{qppcode}
run :: StateT -> Program -> Int -> (StateT, Outcomes, RNG)
run start p seed = evalProg p start (randoms (mkStdGen seed))
\end{qppcode}
The seed makes each run reproducible while keeping evaluation pure.

Measuring one qubit of a Bell pair displays both randomness and collapse:
\begin{qppcode}
peek :: Program
peek = [Unitary bellOp, Measure [0]]

ghci> (s0, o0, _) = evalProg peek (ket [0,0]) choose0
ghci> o0
[False]
ghci> printS s0
1*(ket [0,0])
ghci> (s1, o1, _) = evalProg peek (ket [0,0]) choose1
ghci> o1
[True]
ghci> printS s1
1*(ket [1,1])
\end{qppcode}
Only $00$ and $11$ can occur; seeded runs sample between them reproducibly.

\begin{exercise}[Sampling]
Write
\begin{lstlisting}
freq0 :: StateT -> Program -> Int -> Double
\end{lstlisting}
which runs seeds $1,\ldots,n$ and returns the fraction whose outcome is
\lstinline|[False]|.  Assume $n>0$ and that the program performs one
measurement.  Compare the sampled frequencies of
\begin{lstlisting}
[Unitary H, Measure [0]]
[Unitary H, Unitary H, Measure [0]]
\end{lstlisting}
with their exact Born probabilities.  Explain why two Hadamards do not make
the outcome ``more random''.

One possible outline is a list comprehension with a local pattern:
\begin{lstlisting}
[outcome | seed <- [1..n],
           let (_, outcome, _) = run start program seed]
\end{lstlisting}
Use \lstinline|length| to count matching outcomes and
\lstinline|fromIntegral| before division.
\end{exercise}

\noindent\begin{minipage}{\textwidth}
\section{Which-path information}
\label{sec:whichpath}

Interference requires alternative paths to recombine into the same final
state.  We can prevent this by recording the path in a second qubit.  Insert
a controlled $X$ between the phase shift and the final Hadamard:
\begin{qppcode}
mziWP :: Rational -> Program
mziWP θ =
  [ Unitary ((H ⊗ I) ∘ C X ∘ (R Z θ ⊗ I) ∘ (H ⊗ I))
  , Measure [0]
  ]
\end{qppcode}
\end{minipage}
After \lstinline|C X|, the two alternatives are tagged by orthogonal states
of the second qubit:
\[
  \frac{e^{-i\theta\pi/2}\ket{00}
        +e^{i\theta\pi/2}\ket{11}}{\sqrt2}.
\]
The final $H$ acts only on the first qubit.  It cannot make the $\ket{00}$
and $\ket{11}$ alternatives coincide, so their amplitudes do not interfere.
The first qubit is now measured as $0$ or $1$ with equal probability for
every $\theta$.

No observer needs to read the path record.  Entanglement with two
distinguishable record states is enough.  In a physical machine, uncontrolled
interactions with the environment create such records.  If the record is
inaccessible, the first qubit no longer displays interference.  This loss of
coherence in the subsystem is called \emph{decoherence}.

\begin{exercise}[Erase the path record]
Insert a second \lstinline|C X| immediately after the first.  Predict the
result before sampling it.  Why does the interference fringe return?  State
the conclusion in terms of information, not merely the identity
$\mathrm{CNOT}^2=I$.
\end{exercise}

\section{Reversing a quantum program}
\label{sec:reversal}

A program containing only unitary steps can be collapsed to one unitary
operator and reversed.  The mathematics is
\[
  (UV)^\dagger=V^\dagger U^\dagger,
  \qquad
  (U\otimes V)^\dagger=U^\dagger\otimes V^\dagger,
\]
and, for a controlled operator and a rotation about a Hermitian axis $A$,
\[
  (C(U))^\dagger=C(U^\dagger),\qquad
  R_A(\theta)^\dagger=R_A(-\theta).
\]
The reversal in the first identity is essential: to undo ``first $V$, then
$U$'', undo $U$ first.  Tensor factors do not reverse because they act side
by side rather than in sequence.

QPP exposes the \lstinline|Adjoint| constructor as \lstinline|adj|.  It need
not invert a matrix numerically: recursion over the syntax tree can invert
each gate and reverse composition.  The result remains QPP syntax that any
backend can interpret.  Check the Bell circuit with
\lstinline|printM (evalOp (adj bellOp <> bellOp))|.

\begin{exercise}[Adjoints by structural recursion ($\star$)]
To practise writing a syntax transformation, reimplement this operation.
Prove the composition and tensor identities, then define
\begin{lstlisting}
invert :: QOp -> QOp
\end{lstlisting}
by structural recursion.  Begin with \lstinline|Id|, $X$, $Y$, $Z$, $H$,
tensor, and composition.  Then use the displayed laws for controls and
rotations; an adjoint node cancels an adjoint, while \lstinline|SX| may be
wrapped in \lstinline|Adjoint|.  Compare matrices with \lstinline|adj u|.
The declaration above is abridged, so consult \lstinline|QPP.Syntax| before
claiming an exhaustive definition.  Why can the same method not reverse a
\lstinline|Program| containing \lstinline|Measure|?
\end{exercise}

\begin{samepage}
\section{What to retain}

States are unit vectors up to global phase; unitaries evolve them; measurement
in a chosen orthonormal basis samples the squared magnitudes of the
corresponding amplitudes.  Relative phase produces interference, and
entanglement places information in joint states.  QPP connects these ideas to
typed syntax, structural recursion, interpretation, and effectful syntax:

\begin{center}
\small
\begin{tabular}{p{0.29\textwidth}p{0.31\textwidth}p{0.29\textwidth}}
Quantum concept & QPP representation & Haskell idea \\
\hline
$\ket\psi$, amplitudes & \lstinline|StateT|, \lstinline|ket|, \lstinline|.*|, \lstinline|.+| & values and complex numbers \\
$|\inner{\varphi}{\psi}|^2$ & \lstinline|inner|, \lstinline|prob| & functions and partial application \\
unitary operator & \lstinline|QOp| & algebraic data and syntax trees \\
$U\ket\psi$ & \lstinline|apply (evalOp u) ψ| & syntax interpreted by a backend \\
$U\circ V$, $U\otimes V$, $U^\dagger$ & \lstinline|∘|, \lstinline|⊗|, \lstinline|adj| & recursive program structure \\
measurement & \lstinline|Measure| in a \lstinline|Program| & effectful operations as data \\
repeated experiment & \lstinline|evalProg| with an \lstinline|RNG| & laziness and reproducibility
\end{tabular}
\end{center}

Later lectures build on these correspondences to analyse, transform, and
optimise quantum programs before choosing a backend on which to run them.
\end{samepage}

\end{document}
